\section{Introduction}

Reinforcement learning has become a central component of large language model
(LLM) post-training, particularly for alignment and for reasoning-oriented
tasks with verifiable rewards. A persistent challenge across these settings is
\emph{credit assignment}: trajectories are long, rewards are sparse and
outcome-level, and it is not obvious how a single scalar signal at the end of
a generation should be distributed across hundreds or thousands of token
decisions. Recent work has responded with a rapidly growing toolbox of
token-level credit-assignment methods, including entropy-aware modulation,
process reward models, tree-based prefix values, and reward
redistribution~\citep{cui2025prime,huang2024red,tran2025tempo,wang2025beyond8020,wang2026eapo,chen2026erpo,kazemnejad2024vineppo}.

In parallel, essentially all open-source and academic LLM-RL training uses
parameter-efficient fine-tuning, most commonly LoRA~\citep{hu2021lora}. This
is an empirical reality driven by memory and compute constraints: every
widely used RL framework, from TRL to verl, uses LoRA by default, and recent
results show that LoRA + RL can be dramatically more sample- and
parameter-efficient than full fine-tuning~\citep{tinylora2026,tina2025}. Yet
the two research threads are developed in near-total isolation. Papers on
token-level credit assignment specify methods without reference to the
adaptation strategy, and PEFT is treated as an orthogonal implementation
detail.

In this work we argue that this separation is a mistake, and that the
\emph{interaction} between PEFT and credit assignment is itself a first-class
scientific object. We make the case through three contributions.

\paragraph{Diagnosis: LoRA collapses intrinsic credit signals.} Intrinsic
token-level weighting schemes derive per-token salience from quantities such
as surprisal, entropy reduction, or divergence between the policy and a
reference model. Under LoRA, however, the policy is constrained to a small
low-rank neighbourhood of the reference, and the per-token differences that
these signals measure become approximately uniform across positions. We
formalize this as a collapse of the per-token salience distribution toward
uniform broadcast, characterize it via its Gini coefficient and effective
token count, and show that the collapse reproduces across surprisal-,
entropy-reduction-, and divergence-based weighting. This provides a mechanistic
explanation for the empirically observed failure of full-token
GRPO~+~LoRA~\citep{tokenefficientrl2025}: the weighting signal itself is
degraded before training ever begins.

\paragraph{Method: Adapter-Residual Token Weighting.} Rather than attempting
to recover per-token structure from signals that are intrinsically flattened
under LoRA, we propose to measure per-token salience directly from the
adapter's own contribution to the forward pass. Adapter-Residual Token
Weighting (ARTW) sets the salience at position $t$ equal to the norm of the
adapter residual, $\|h^{\text{adapted}}_t - h^{\text{base}}_t\|_2$, computed
via a forward pass with the adapter disabled. This signal is positive
whenever the adapter is active, non-uniform across positions by construction
(because adapter input activations are non-uniform), and requires no
additional networks, reward models, or tree constructions. It reuses exactly
the adapter-disabled forward pass already needed for KL regularization and
divergence weighting.

\paragraph{Framework: Unified diagnostic evaluation.} We reframe our
experimental comparison from "which weighting scheme wins" to "which
weighting schemes remain non-degenerate under PEFT." We report a unified set
of concentration diagnostics (Gini coefficient, effective token count,
greedy accuracy, pass@$k$) across uniform, surprisal, entropy-reduction,
divergence, critic, and ARTW weighting under both RLOO and GRPO, and
characterize how these diagnostics change with LoRA rank. The framework
makes it possible to distinguish between credit-assignment methods that
\emph{could} work in principle and methods that actually
\emph{retain signal} under the adaptation strategy in use.

The remainder of the paper develops these contributions. Section
\ref{sec:related-work} positions the work relative to the token-level
credit-assignment and PEFT-for-RL literatures. Section \ref{sec:methods}
presents the weighting schemes, proves the collapse result formally, and
defines ARTW. Section \ref{sec:theoretical} analyses the bias--variance
behaviour of the method. Section \ref{sec:experiments} reports the diagnostic
and performance comparison on GSM8K with Qwen2.5 under multiple LoRA ranks.
